\documentclass[a4paper,11pt,reqno]{amsart}

% =========================================================================
%  Instrumentación Electrónica — PreLab de la sesión 1 (amplificadores operacionales)
%  Documento autocontenido para pasar a manuscrito. Compilar con: latexmk -pdf main.tex
%  Ecuaciones numeradas de (1) en adelante, a la derecha.
% =========================================================================
% =========================================================================
%  Instrumentación Electrónica — PreLab de la sesión 1 (amplificadores operacionales)
%  Estilo heredado de los apuntes de Electromagnetismo I (Temas 0 y 1):
%  amsart, Palatino, versalitas, enunciados numerados por sección, paleta de
%  azules y grises sobrios. Ecuaciones numeradas a la derecha (opción reqno).
% =========================================================================

% --- IDIOMA Y CODIFICACIÓN ------------------------------------------------
\usepackage[utf8]{inputenc}
\usepackage[T1]{fontenc}
\usepackage[spanish,es-tabla,es-nodecimaldot]{babel}

% --- TIPOGRAFÍA (Palatino) --------------------------------------------------
\usepackage{mathpazo}
\linespread{1.05}
\emergencystretch=2em
\usepackage[final]{microtype}

% --- MATEMÁTICAS -------------------------------------------------------------
\usepackage{amsmath, amssymb, mathtools}
\usepackage{bm}

% --- UNIDADES (coma decimal, como en el resto de los apuntes) ---------------
\usepackage{siunitx}
\sisetup{output-decimal-marker={,}, group-separator={\,}, group-minimum-digits=5,
         exponent-product=\cdot, per-mode=symbol, range-phrase={--},
         list-final-separator={ y }, list-pair-separator={ y }}

% --- PÁGINA, COLORES Y GRÁFICOS --------------------------------------------
\usepackage[a4paper,margin=2.6cm]{geometry}
\usepackage{graphicx}
\usepackage{xcolor}
% Paleta base (idéntica a la de los apuntes de Electromagnetismo I)
\definecolor{azulprof}{HTML}{1F3A5F}    % azul oscuro
\definecolor{azulacero}{HTML}{3B6EA5}   % azul medio
\definecolor{azulclaro}{HTML}{8FB3D9}   % azul claro
\definecolor{azulfondo}{HTML}{F3F7FB}   % fondo tenue de definiciones
\definecolor{grispizarra}{HTML}{6B7280} % gris
\definecolor{grisclaro}{HTML}{D4D8DE}   % gris claro
\definecolor{acento}{HTML}{B8800A}      % acento: ámbar oscuro
\definecolor{rojoteo}{HTML}{B8800A}     % señales de salida, resultados
\definecolor{verdeteo}{HTML}{2F6F68}    % corrientes y sentidos de referencia
\definecolor{opfondo}{HTML}{E4EEF8}     % relleno del triángulo del AO

\usepackage{tikz}
\usetikzlibrary{babel, arrows.meta, calc, positioning, decorations.markings,
                patterns, intersections}
\usepackage[european, straightvoltages, RPvoltages]{circuitikz}
\usepackage{pgfplots}
\pgfplotsset{compat=1.17}
\usepgfplotslibrary{fillbetween}
\usepackage{booktabs}
\usepackage{needspace}
\usepackage{array}
\usepackage{multirow}
\usepackage[inline]{enumitem}
\setlist[itemize]{leftmargin=1.6em, itemsep=2pt, topsep=4pt}
\setlist[enumerate]{leftmargin=1.8em, itemsep=2pt, topsep=4pt}

% --- ENLACES Y REFERENCIAS CRUZADAS ----------------------------------------
\usepackage[colorlinks=true, linkcolor=azulacero, citecolor=azulacero,
            urlcolor=azulacero, plainpages=false, pdfusetitle=false,
            pdftitle={PreLab. Sesión 1: amplificadores operacionales},
            pdfauthor={Daniel Martín Guerrero}]{hyperref}
\usepackage[spanish,nameinlink]{cleveref}

% --- CABECERAS ---------------------------------------------------------------
\pagestyle{headings}
\makeatletter
\def\ps@headings{%
  \let\@mkboth\markboth
  \def\@evenhead{\small\scshape\normalfont\thepage\hfil Instrumentación Electrónica}%
  \def\@oddhead{\small\scshape\normalfont PreLab. Sesión 1: amplificadores operacionales\hfil\thepage}%
  \let\@evenfoot\@empty \let\@oddfoot\@empty}
\makeatother

% Subsecciones como títulos en bloque (no «run-in»), con aire antes y después
\usepackage{titlesec}
\titleformat{\subsection}[block]{\normalfont\bfseries\color{azulprof}}{\thesubsection.}{0.6em}{}
\titlespacing*{\subsection}{0pt}{3ex plus 1ex minus .5ex}{1.6ex plus .4ex}
\setcounter{tocdepth}{2}
\makeatletter
\def\@@and{y}
\renewcommand{\keywordsname}{Palabras clave}
% Índice: número de sección en una caja de anchura fija (evita el solape de amsart)
\renewcommand{\tocsection}[3]{\indentlabel{\@ifnotempty{#2}{\makebox[1.6em][l]{#2.}}}#3}
\renewcommand{\tocsubsection}[3]{\indentlabel{\@ifnotempty{#2}{\makebox[2.4em][l]{#2.}}}#3}
\def\l@section{\@tocline{1}{2pt}{0pt}{1.6em}{}}
\def\l@subsection{\@tocline{2}{0pt}{1.6em}{2.4em}{}}
\def\l@part{\@tocline{-1}{8pt plus 2pt}{0pt}{0pt}{\bfseries\color{azulprof}}}
\makeatother

% =========================================================================
%               NOTACIÓN Y MACROS DE USO FRECUENTE
% =========================================================================
\newcommand{\dif}{\mathrm{d}}
\newcommand{\eul}{\mathrm{e}}
\newcommand{\jj}{\mathrm{j}}                         % unidad imaginaria (j)
\newcommand{\R}{\mathbb{R}}
\newcommand{\C}{\mathbb{C}}
\newcommand{\abs}[1]{\lvert #1\rvert}
\newcommand{\defeq}{\mathrel{\overset{\mathrm{def}}{=}}}
\newcommand{\Lt}{\mathcal{L}}                        % transformada de Laplace
\DeclareMathOperator{\sen}{sen}
\DeclareMathOperator{\sgn}{sgn}
\DeclareMathOperator{\arctg}{arctg}
\DeclareMathOperator{\Real}{Re}
\renewcommand{\sin}{\sen}                            % «sen» en todo el texto
% Magnitudes de circuito
\newcommand{\vin}{v_{\mathrm{in}}}
\newcommand{\vout}{v_{\mathrm{out}}}
\newcommand{\Vin}{V_{\mathrm{in}}}
\newcommand{\Vout}{V_{\mathrm{out}}}
\newcommand{\Vsat}{V_{\mathrm{sat}}}
\newcommand{\Vcc}{V_{\mathrm{CC}}}
\newcommand{\Vee}{V_{\mathrm{EE}}}
\newcommand{\Rin}{R_{\mathrm{in}}}
\newcommand{\Rout}{R_{\mathrm{out}}}
\newcommand{\Zin}{Z_{\mathrm{in}}}
\newcommand{\Zout}{Z_{\mathrm{out}}}
\newcommand{\Av}{A_{V}}
\newcommand{\vp}{v_{+}}
\newcommand{\vm}{v_{-}}
\newcommand{\dB}{\,\mathrm{dB}}
\newcommand{\pend}[1]{\textcolor{grispizarra}{\colorbox{grisclaro}{\strut\,#1\,}}}

% =========================================================================
%        ENTORNOS FORMALES (enunciados numerados, estilo artículo)
% =========================================================================
\usepackage{amsthm}
\usepackage[most]{tcolorbox}

\newtheoremstyle{cienplano}{10pt}{10pt}{\itshape}{}%
  {\bfseries\color{azulprof}}{.}{0.6em}%
  {\thmname{#1}\thmnumber{ #2}\thmnote{\normalfont\upshape\ (#3)}}
\newtheoremstyle{ciendef}{10pt}{10pt}{\normalfont}{}%
  {\bfseries\color{azulprof}}{.}{0.6em}%
  {\thmname{#1}\thmnumber{ #2}\thmnote{\normalfont\upshape\ (#3)}}
\newtheoremstyle{cienobs}{8pt}{8pt}{\normalfont}{}%
  {\itshape\color{azulprof}}{.}{0.6em}%
  {\thmname{#1}\thmnumber{ #2}\thmnote{\normalfont\upshape\ (#3)}}

\theoremstyle{cienplano}
\newtheorem{proposicion}{Proposición}
\newtheorem{teorema}[proposicion]{Teorema}
\newtheorem{corolario}[proposicion]{Corolario}
\theoremstyle{ciendef}
\newtheorem{ejemplo}[proposicion]{Ejemplo}
\newtheorem{ejercicio}[proposicion]{PreLab}
\newtheorem{hipotesis}[proposicion]{Hipótesis}
% Definiciones: único entorno enmarcado (caja fina, fondo tenue, título en el borde)
\newcounter{definicion}
\renewcommand{\thedefinicion}{\theproposicion}
\newtcolorbox{definicionbox}[1]{%
  enhanced, breakable, sharp corners,
  colback=azulfondo, colframe=azulclaro, boxrule=0.4pt,
  left=9pt, right=9pt, top=5pt, bottom=5pt,
  before skip=12pt, after skip=12pt,
  attach boxed title to top left={yshift*=-\tcboxedtitleheight/2, xshift=8pt},
  boxed title style={colback=white, colframe=white, sharp corners, boxrule=0pt,
                     left=3pt, right=3pt, top=1pt, bottom=1pt},
  fonttitle=\bfseries\color{azulprof}, coltitle=azulprof,
  title={Definición~\thedefinicion\ifx\relax#1\relax\else\ (#1)\fi}}
\newenvironment{definicion}[1][]
  {\par\stepcounter{proposicion}\refstepcounter{definicion}\begin{definicionbox}{#1}}
  {\end{definicionbox}}

% Enunciado de cada recuadro del PreLab: misma caja, con filete ámbar (los
% recuadros amarillos del guion), para distinguir lo pedido de lo resuelto.
\newcounter{prelab}
\newtcolorbox{prelabbox}[1]{%
  enhanced, sharp corners,
  colback=white, colframe=acento!70, boxrule=0.5pt, leftrule=2.2pt,
  left=9pt, right=9pt, top=5pt, bottom=5pt,
  before skip=12pt, after skip=10pt,
  attach boxed title to top left={yshift*=-\tcboxedtitleheight/2, xshift=8pt},
  boxed title style={colback=white, colframe=white, sharp corners, boxrule=0pt,
                     left=3pt, right=3pt, top=1pt, bottom=1pt},
  fonttitle=\bfseries\color{acento!85!black}, coltitle=acento!85!black,
  title={PreLab~\theprelab\ifx\relax#1\relax\else\ (#1)\fi}}
\newenvironment{prelab}[1][]
  {\par\refstepcounter{prelab}\begin{prelabbox}{#1}}
  {\end{prelabbox}}

\theoremstyle{cienobs}
\newtheorem{observacion}[proposicion]{Observación}

% Resolución: misma maquinaria que «proof», con otro rótulo
\renewcommand{\qedsymbol}{\ensuremath{\blacksquare}}
\newenvironment{resolucion}[1][Resolución]
  {\begin{proof}[\itshape\color{azulprof}#1]}{\end{proof}}

% Resultado final enmarcado (como \boxed, pero con el filete de la paleta)
\newcommand{\resultado}[1]{\fcolorbox{azulacero}{white}{\ensuremath{\displaystyle #1}}}

\crefname{proposicion}{proposición}{proposiciones}
\Crefname{proposicion}{Proposición}{Proposiciones}
\crefname{teorema}{teorema}{teoremas}
\Crefname{teorema}{Teorema}{Teoremas}
\crefname{corolario}{corolario}{corolarios}
\Crefname{corolario}{Corolario}{Corolarios}
\crefname{definicion}{definición}{definiciones}
\Crefname{definicion}{Definición}{Definiciones}
\crefname{ejemplo}{ejemplo}{ejemplos}
\Crefname{ejemplo}{Ejemplo}{Ejemplos}
\crefname{hipotesis}{hipótesis}{hipótesis}
\Crefname{hipotesis}{Hipótesis}{Hipótesis}
\crefname{prelab}{PreLab}{PreLabs}
\Crefname{prelab}{PreLab}{PreLabs}
\crefname{observacion}{observación}{observaciones}
\Crefname{observacion}{Observación}{Observaciones}
\crefname{section}{sección}{secciones}
\Crefname{section}{Sección}{Secciones}
\crefname{subsection}{sección}{secciones}
\Crefname{subsection}{Sección}{Secciones}
\crefname{table}{tabla}{tablas}
\Crefname{table}{Tabla}{Tablas}
\crefname{figure}{figura}{figuras}
\Crefname{figure}{Figura}{Figuras}
\crefname{equation}{ecuación}{ecuaciones}
\Crefname{equation}{Ecuación}{Ecuaciones}
\crefformat{equation}{ecuación~(#2#1#3)}
\Crefformat{equation}{Ecuación~(#2#1#3)}
\crefrangeformat{equation}{ecuaciones~(#3#1#4)--(#5#2#6)}
\crefmultiformat{equation}{ecuaciones~(#2#1#3)}{ y~(#2#1#3)}{, (#2#1#3)}{ y~(#2#1#3)}

\usepackage[font=small, labelfont={sc,color=azulprof}, labelsep=period,
            margin=0.6cm, justification=justified]{caption}
\usepackage{subcaption}

% Referencia a un panel de figura: «figura 5\,(b)», enlazada en bloque
\newcommand{\figpanel}[2]{\hyperref[#1]{figura~\ref*{#1}\,\textup{(#2)}}}
\newcommand{\Figpanel}[2]{\hyperref[#1]{Figura~\ref*{#1}\,\textup{(#2)}}}

% =========================================================================
%        ESTILOS GRÁFICOS COMUNES (TikZ / circuitikz / pgfplots)
% =========================================================================
% =========================================================================
%  Estilos gráficos comunes: TikZ, circuitikz y pgfplots
% =========================================================================

% --- Estilos TikZ generales ------------------------------------------------
\tikzset{
  >=Stealth,
  eje/.style={->, line width=0.6pt, grispizarra},
  fino/.style={line width=0.4pt, grispizarra},
  punteado/.style={densely dotted, line width=0.5pt, grispizarra},
  etiq/.style={font=\small, text=black!85},
  nota/.style={font=\scriptsize, text=grispizarra, align=center},
  % corriente dibujada junto a un hilo (flecha verde con rótulo)
  corr/.style={-{Stealth[length=4.5pt,width=4pt]}, line width=0.8pt, verdeteo},
  % nudo con tensión anotada
  tens/.style={font=\small, text=azulprof},
}

% --- circuitikz: trazo y colores sobrios -----------------------------------
\ctikzset{
  bipoles/thickness=1,
  bipoles/length=1.15cm,
  tripoles/op amp/width=1.5,
  tripoles/op amp/height=1.35,
  tripoles/op amp/font=\small,
  current/distance=0.45,
  label/align=smart,
}
\tikzset{
  circuito/.style={line width=0.55pt, color=black!85, font=\small},
  ao/.style={fill=opfondo},
}

% --- Estilo base de ejes pgfplots ------------------------------------------
\pgfplotsset{
  tema/.style={
    tick label style={font=\scriptsize}, label style={font=\small},
    axis lines=left, axis line style={grispizarra},
    grid=major, grid style={grisclaro, line width=0.4pt},
    every tick/.style={grispizarra},
    legend style={font=\scriptsize, draw=none, fill=white, fill opacity=0.85,
                  text opacity=1},
    legend cell align=left,
    /pgf/number format/use comma,
  },
  % ejes centrados (formas de onda y características de transferencia)
  temacentro/.style={
    tick label style={font=\scriptsize}, label style={font=\small},
    axis lines=middle, axis line style={grispizarra, -{Stealth[length=4pt]}},
    every tick/.style={grispizarra},
    legend style={font=\scriptsize, draw=none, fill=white, fill opacity=0.85,
                  text opacity=1},
    legend cell align=left,
    /pgf/number format/use comma,
  },
  senal entrada/.style={azulacero, line width=1.0pt},
  senal salida/.style={rojoteo, line width=1.1pt},
  asintota/.style={grispizarra, densely dashed, line width=0.6pt},
}



% =========================================================================
%                     VARIABLES PARA LA PORTADA
% =========================================================================
\newcommand{\Universidad}{Universidad de Salamanca}
\newcommand{\Facultad}{Facultad de Ciencias}
\newcommand{\Titulacion}{Grado en Física}
\newcommand{\Asignatura}{Instrumentación Electrónica}
\newcommand{\NumTema}{Laboratorio de circuitos, sesión 1}
\newcommand{\TituloTema}{PreLab:\\amplificadores operacionales}
\newcommand{\tipotrabajo}{PreLab. Cálculos de diseño}
\newcommand{\miNombre}{Daniel Martín Guerrero}
\newcommand{\miRotuloProfesor}{Profesores}
\newcommand{\miProfesor}{Yahya Moubarak Meziani\\María Susana Pérez Santos\\Jesús Enrique Velázquez Pérez}
\newcommand{\miDetalleProfesor}{Área de Electrónica}
\newcommand{\SubtituloTema}{Amplificador inversor, amplificador no inversor\\y amplificador inversor generalizado como filtro paso bajo}
\newcommand{\miUbicacion}{Salamanca, curso 2026/2027}

\title{PreLab de la sesión 1: amplificadores operacionales}
\author{Daniel Martín Guerrero}
\address{Grado en Física, Facultad de Ciencias, Universidad de Salamanca,
         Salamanca, España}
\date{}

\begin{document}
\pagenumbering{alph}
%%%%%%%%%%%%%%%%%%%%%%%%%%%%%%%%%%%%%%%%%%%%%%%%%%%%%%%%%%%%%%%%%%%%%%%%
% Portada formal (estilo tesis): logotipos, identificación institucional,
% título y datos de autoría. Banda vertical a la derecha, como en el
% «Cálculo integral multivariable» de L. Navas. Requiere en main.tex las
% macros \Universidad, \Facultad, \Titulacion, \Asignatura, \NumTema,
% \TituloTema, \SubtituloTema, \tipotrabajo, \miNombre, \miProfesor,
% \miUbicacion, \miRotuloProfesor, \miDetalleProfesor.
%%%%%%%%%%%%%%%%%%%%%%%%%%%%%%%%%%%%%%%%%%%%%%%%%%%%%%%%%%%%%%%%%%%%%%%%

\begin{titlepage}
\thispagestyle{empty}
\pagecolor{white}

\begin{tikzpicture}[remember picture, overlay]
  \coordinate (O) at (current page.south west);

  % Velo muy tenue en la mitad superior (cabecera institucional)
  \shade[top color=azulfondo, bottom color=white]
    ($(O)+(0,18.5cm)$) rectangle ($(O)+(21cm,29.7cm)$);

  % Banda vertical azul, a la derecha del bloque de texto
  \fill[azulacero] ($(O)+(19.0cm,2.6cm)$) rectangle ($(O)+(19.14cm,27.4cm)$);

  % Logotipos
  \node[anchor=north west, inner sep=0] at ($(O)+(2.4cm,27.7cm)$)
    {\includegraphics[height=1.55cm]{logo_usal.png}};
  \node[anchor=north east, inner sep=0] at ($(O)+(17.4cm,27.7cm)$)
    {\includegraphics[height=1.55cm]{logo_facultad.png}};
  \draw[grispizarra!60, line width=0.4pt] ($(O)+(2.4cm,25.7cm)$) -- ($(O)+(17.4cm,25.7cm)$);

  % Identificación institucional
  \node[anchor=north, align=center, inner sep=0, text width=15cm] at ($(O)+(9.9cm,25.2cm)$)
    {{\color{azulprof}\scshape\Large\textls[120]{\Universidad}}\\[5pt]
     {\color{grispizarra}\scshape\large\textls[80]{\Facultad}}\\[3pt]
     {\color{grispizarra}\scshape\large\textls[80]{\Titulacion}}};

  % Asignatura y título
  \node[anchor=north, align=center, inner sep=0, text width=15cm] at ($(O)+(9.9cm,19.2cm)$)
    {{\color{azulacero}\scshape\Large\textls[160]{\Asignatura}}\\[4pt]
     {\color{grispizarra}\scshape\large\textls[100]{\NumTema}}\\[22pt]
     {\color{azulprof}\fontsize{30}{36}\selectfont\bfseries\TituloTema}\\[14pt]
     {\color{azulacero}\rule{3.2cm}{0.8pt}}\\[14pt]
     {\color{black!80}\itshape\large\SubtituloTema}};

  % Epígrafe: los tres resultados del PreLab
  \node[anchor=north, align=center, inner sep=0, text width=15cm] at ($(O)+(9.9cm,11.0cm)$)
    {\color{grispizarra}\small
     $\displaystyle A_{V}=-\frac{R_{2}}{R_{1}}\qquad
      A_{V}=1+\frac{R_{2}}{R_{1}}\qquad
      H(s)=-\frac{Z_{2}(s)}{Z_{1}(s)}$};

  % Autoría y dirección
  \draw[grispizarra!60, line width=0.4pt] ($(O)+(2.4cm,8.6cm)$) -- ($(O)+(17.4cm,8.6cm)$);
  \node[anchor=north west, inner sep=0, text width=7cm, align=left] at ($(O)+(2.4cm,8.2cm)$)
    {{\scshape\small\color{azulacero}\textls[80]{Autor}}\\[4pt]
     {\large\color{black!85}\miNombre}\\[2pt]
     {\small\color{grispizarra}Grado en Física}};
  \node[anchor=north west, inner sep=0, text width=7cm, align=left] at ($(O)+(10.4cm,8.2cm)$)
    {{\scshape\small\color{azulacero}\textls[80]{\miRotuloProfesor}}\\[4pt]
     {\color{black!85}\miProfesor}\\[3pt]
     {\small\color{grispizarra}\miDetalleProfesor}};

  % Tipo de documento, lugar y curso
  \node[anchor=south, align=center, inner sep=0, text width=15cm] at ($(O)+(9.9cm,2.9cm)$)
    {{\color{azulprof}\scshape\large\textls[100]{\tipotrabajo}}\\[4pt]
     {\color{grispizarra}\scshape\small\textls[80]{\miUbicacion}}};
\end{tikzpicture}

\end{titlepage}

\clearpage
\pagenumbering{arabic}
\setcounter{page}{1}

\maketitle

% =========================================================================
\section{Introducción}\label{sec:introduccion}
% =========================================================================
En este PreLab se obtienen, con el modelo ideal del amplificador operacional (AO), las ganancias de
tres circuitos realimentados: el amplificador inversor ($A_V=-10$), el no inversor ($A_V=+11$) y el
inversor generalizado, que con un condensador en la realimentación se comporta como un filtro paso bajo.

Un AO amplifica la diferencia entre sus entradas no inversora ($+$) e inversora ($-$),
$\vout=A\,(\vp-\vm)$, con una ganancia en lazo abierto $A$ muy grande; tiene una resistencia de entrada
$\Rin$ entre ambas entradas y una resistencia de salida $\Rout$. El \emph{AO ideal} cumple
\[
  \Rin\to\infty,\qquad \Rout\to0,\qquad A\to\infty .
\]
Si trabaja en zona lineal gracias a una \emph{realimentación negativa} (una red que lleva parte de la
salida a la entrada inversora), se derivan las dos reglas que usaremos en todos los cálculos:
\begin{itemize}[leftmargin=3em]
  \item[(R1)] \emph{corrientes de entrada nulas}, $i_+=i_-=0$, porque entre las entradas solo hay $\Rin\to\infty$;
  \item[(R2)] \emph{cortocircuito virtual}, $\vp=\vm$, porque la salida es finita ($\abs{\vout}\le\Vsat$) y
        $\vp-\vm=\vout/A\to0$ cuando $A\to\infty$.
\end{itemize}
Además, con $\Rout=0$ la tensión de salida no depende de la carga. Cada circuito se resuelve imponiendo
(R1) y (R2) y aplicando la ley de Kirchhoff de corrientes en el nudo inversor y la ley de Ohm.

% =========================================================================
\section{Circuito 1: amplificador inversor}\label{sec:inversor}
% =========================================================================
\begin{prelab}[Amplificador inversor]
Usando el modelo de amplificador operacional ideal visto en teoría, encuéntrese la ganancia en
tensión $A_V=\Vout/\Vin$ en función de las resistencias del circuito de la \cref{fig:inversor}. Si
imponemos $R_1=1\,\mathrm{k}\Omega$, determínese el valor de $R_2$ para conseguir $A_V=-10$ a $1\,\mathrm{kHz}$.
\end{prelab}

% Figura: amplificador inversor (circuito 1) con tensiones y corrientes del análisis.
\begin{figure}[ht]
\centering
\begin{circuitikz}[circuito]
  \draw (0,0) node[op amp, ao] (A) {};
  % nudo inversor y red de realimentación
  \draw (A.-) -- ++(-0.6,0) coordinate (N);
  \draw (N) to[R, l_=$R_1$, i<^=$i_1$] ++(-2.4,0) coordinate (E);
  \draw (N) |- ++(0,1.5) coordinate (T);
  \draw (A.out) -- ++(0.7,0) coordinate (O);
  \draw (T) to[R, l=$R_2$, i>^=$i_2$] (T -| O) -- (O);
  % generador
  \draw (E) to[short, -o] ++(-0.3,0) coordinate (Ein) node[above, tens] {$\vin$};
  \draw (Ein) -- ++(-0.9,0) coordinate (G) to[sV, l_=$v_g$] ++(0,-2.0) node[ground] {};
  % entrada no inversora a tierra
  \draw (A.+) -- ++(-0.6,0) coordinate (P) -- ++(0,-0.45) node[ground] {};
  % carga y salida
  \draw (O) to[short, *-o] ++(1.0,0) coordinate (S) node[right, tens] {$\vout$};
  \draw (O) ++(0.5,0) to[R, l=$R_L$] ++(0,-2.0) node[ground] {};
  \fill (N) circle (1.6pt);
  % anotaciones del análisis
  \node[tens, anchor=east, font=\scriptsize] at ($(N)+(-0.08,0.75)$) {$\vm=0$};
  \draw[azulprof, line width=0.3pt] ($(N)+(-0.08,0.72)$) -- ($(N)+(-0.02,0.06)$);
  \node[tens, anchor=east, font=\scriptsize] at ($(P)+(-0.1,-0.25)$) {$\vp=0$};
\end{circuitikz}
\caption{Amplificador inversor (figura~2 del guion), con $R_1=1\,\mathrm{k}\Omega$,
$R_L=10\,\mathrm{k}\Omega$ y la alimentación omitida. Se indican las magnitudes del análisis: el nudo
inversor es una tierra virtual ($\vm=\vp=0$) y no absorbe corriente, de modo que $i_1=i_2$.}
\label{fig:inversor}
\end{figure}


\begin{resolucion}
AO ideal con realimentación negativa a través de $R_2$; sentidos de corriente de la \cref{fig:inversor}.

\textbf{1. Tierra virtual.} La entrada $+$ está a tierra; por (R2),
\begin{equation}\label{eq:i-vm}
  \vm=\vp=0 .
\end{equation}

\textbf{2. Ley de Ohm en $R_1$}, entre $\vin$ y $\vm$:
\begin{equation}\label{eq:i-i1}
  i_1=\frac{\vin-\vm}{R_1}=\frac{\vin}{R_1} .
\end{equation}

\textbf{3. Kirchhoff en el nudo inversor}, con $i_-=0$ por (R1):
\begin{equation}\label{eq:i-i2}
  i_1=i_2+i_-\quad\Longrightarrow\quad i_2=i_1=\frac{\vin}{R_1} .
\end{equation}

\textbf{4. Ley de Ohm en $R_2$}, recorrida por $i_2$ del nudo inversor a la salida:
\begin{equation}\label{eq:i-vout}
  \vm-\vout=R_2\,i_2\quad\Longrightarrow\quad \vout=-\frac{R_2}{R_1}\,\vin .
\end{equation}

\textbf{Resultado.}
\begin{equation}\label{eq:i-Av}
  \resultado{A_V=\frac{\vout}{\vin}=-\frac{R_2}{R_1}}
\end{equation}

\textbf{Diseño.} Con $A_V=-10$ y $R_1=1\,\mathrm{k}\Omega$ en \eqref{eq:i-Av}:
\begin{equation}\label{eq:i-R2}
  R_2=-A_V\,R_1=10\cdot1\,\mathrm{k}\Omega\quad\Longrightarrow\quad\resultado{R_2=10\,\mathrm{k}\Omega}
\end{equation}

El signo menos indica contrafase ($180^{\circ}$). La ganancia no depende de la frecuencia, de modo que el
diseño vale a $1\,\mathrm{kHz}$, ni de $A$ ni de $R_L$: solo del cociente $R_2/R_1$.
\end{resolucion}

% =========================================================================
\section{Circuito 2: amplificador no inversor}\label{sec:no-inversor}
% =========================================================================
\begin{prelab}[Amplificador no inversor]
Usando un modelo de amplificador operacional ideal, encuéntrese la ganancia en tensión
$A_V=\Vout/\Vin$ en función de las resistencias del circuito de la \cref{fig:no-inversor}. Si imponemos
$R_1=1\,\mathrm{k}\Omega$, determínese el valor de $R_2$ para conseguir $A_V=+11$ a $1\,\mathrm{kHz}$.
\end{prelab}

% Figura: amplificador no inversor (circuito 2) con tensiones y corrientes del análisis.
\begin{figure}[ht]
\centering
\begin{circuitikz}[circuito]
  \draw (0,0) node[op amp, ao, noinv input up] (A) {};
  % nudo inversor (abajo) y red de realimentación
  \draw (A.out) -- ++(0.7,0) coordinate (O);
  \draw (A.-) -- ++(-0.6,0) coordinate (N) -- ++(0,-1.0) coordinate (M);
  \draw (M) -- ++(-1.0,0) coordinate (G1) to[R, l_=$R_1$, i>^=$i_1$] ++(0,-1.7) node[ground] {};
  \draw (M) to[R, l_=$R_2$, i<_=$i_2$] (M -| O) -- (O);
  % generador en la entrada no inversora
  \draw (A.+) to[short, -o] ++(-2.0,0) coordinate (Ein) node[above, tens] {$\vin$};
  \draw (Ein) -- ++(-0.8,0) to[sV, l_=$v_g$] ++(0,-2.0) node[ground] {};
  % carga y salida
  \draw (O) to[short, *-o] ++(1.0,0) node[right, tens] {$\vout$};
  \draw (O) ++(0.5,0) to[R, l=$R_L$] ++(0,-2.0) node[ground] {};
  \fill (M) circle (1.6pt);
  \node[tens, anchor=east, font=\scriptsize] at ($(N)!0.5!(M)+(-0.08,0)$) {$\vm=\vin$};
\end{circuitikz}
\caption{Amplificador no inversor (figura~3 del guion), con $R_1=1\,\mathrm{k}\Omega$,
$R_L=10\,\mathrm{k}\Omega$ y la alimentación omitida. La entrada $+$ no absorbe corriente; el nudo
inversor sigue a la entrada ($\vm=\vin$) y $R_1$--$R_2$ forman un divisor de $\vout$.}
\label{fig:no-inversor}
\end{figure}


\begin{resolucion}
AO ideal con realimentación negativa a través de $R_2$; sentidos de corriente de la \cref{fig:no-inversor}.

\textbf{1. Cortocircuito virtual.} La señal entra por $+$ sin caída de tensión ($i_+=0$, R1); por (R2),
\begin{equation}\label{eq:n-vm}
  \vm=\vp=\vin .
\end{equation}

\textbf{2. Ley de Ohm en $R_1$}, entre $\vm$ y tierra:
\begin{equation}\label{eq:n-i1}
  i_1=\frac{\vm-0}{R_1}=\frac{\vin}{R_1} .
\end{equation}

\textbf{3. Kirchhoff en el nudo inversor}, con $i_-=0$ por (R1):
\begin{equation}\label{eq:n-i2}
  i_2=i_1+i_-\quad\Longrightarrow\quad i_2=i_1=\frac{\vin}{R_1} .
\end{equation}

\textbf{4. Ley de Ohm en $R_2$}, recorrida por $i_2$ de la salida al nudo inversor:
\begin{equation}\label{eq:n-vout}
  \vout-\vm=R_2\,i_2\quad\Longrightarrow\quad \vout=\vin+\frac{R_2}{R_1}\,\vin=\Bigl(1+\frac{R_2}{R_1}\Bigr)\vin .
\end{equation}

\textbf{Resultado.}
\begin{equation}\label{eq:n-Av}
  \resultado{A_V=\frac{\vout}{\vin}=1+\frac{R_2}{R_1}}
\end{equation}

\textbf{Diseño.} Con $A_V=11$ y $R_1=1\,\mathrm{k}\Omega$ en \eqref{eq:n-Av}:
\begin{equation}\label{eq:n-R2}
  R_2=(A_V-1)\,R_1=10\cdot1\,\mathrm{k}\Omega\quad\Longrightarrow\quad\resultado{R_2=10\,\mathrm{k}\Omega}
\end{equation}

$A_V>0$: salida en fase con la entrada; además $A_V\ge1$ siempre. Como en el inversor, la ganancia no
depende de la frecuencia, ni de $A$ ni de $R_L$.
\end{resolucion}

% =========================================================================
\section{Circuito 3: amplificador inversor generalizado y filtro paso bajo}\label{sec:filtro}
% =========================================================================
\begin{prelab}[Inversor generalizado y filtro]
Sea el circuito de la \figpanel{fig:generalizado}{a}, en el que las impedancias $Z_1$ y $Z_2$ pueden ser no
resistivas y, por lo tanto, tener dependencia en frecuencia: obténgase la ganancia en voltaje del circuito.
Esta ganancia será función de la frecuencia salvo que ambas impedancias sean puramente resistivas.
Considérese el caso de que $Z_1=R$ sea real y $Z_2$ una composición en paralelo de una resistencia $R_F$ y
un condensador $C$ (\figpanel{fig:generalizado}{b}). Demuéstrese que la dependencia de la función de
transferencia con la frecuencia $f$ es la de un filtro paso bajo: para frecuencias suficientemente bajas
la ganancia en tensión es constante, mientras que para frecuencias suficientemente altas disminuye.
\end{prelab}

% Figura: (a) amplificador inversor generalizado; (b) caso Z1=R, Z2=RF||C (filtro paso bajo).
\begin{figure}[ht]
\centering
\begin{minipage}[b]{0.46\textwidth}\centering
\begin{circuitikz}[circuito, scale=0.88, transform shape]
  \draw (0,0) node[op amp, ao] (A) {};
  \draw (A.-) -- ++(-0.5,0) coordinate (N);
  \draw (N) to[generic, l_=$Z_1$, i<^=$I_1$] ++(-2.2,0) to[short, -o] ++(-0.3,0) node[left, tens] {$\Vin$};
  \draw (A.out) -- ++(0.6,0) coordinate (O);
  \draw (N) -- ++(0,1.4) coordinate (T) to[generic, l=$Z_2$, i>^=$I_2$] (T -| O) -- (O);
  \draw (O) to[short, *-o] ++(0.6,0) node[right, tens] {$\Vout$};
  \draw (A.+) -- ++(-0.5,0) -- ++(0,-0.4) node[ground] {};
  \fill (N) circle (1.6pt);
\end{circuitikz}\\[2pt]
{\small (a)}
\end{minipage}\hfill
\begin{minipage}[b]{0.50\textwidth}\centering
\begin{circuitikz}[circuito, scale=0.88, transform shape]
  \draw (0,0) node[op amp, ao] (A) {};
  \draw (A.-) -- ++(-0.5,0) coordinate (N);
  \draw (N) to[R, l_=$R$] ++(-2.2,0) to[short, -o] ++(-0.3,0) node[left, tens] {$\Vin$};
  \draw (A.out) -- ++(0.6,0) coordinate (O);
  \draw (N) -- ++(0,1.2) coordinate (T1) to[C, l_=$C$] (T1 -| O) -- (O);
  \draw (T1) -- ++(0,1.0) coordinate (T2) to[R, l=$R_F$] (T2 -| O) -- (T1 -| O);
  \fill (T1) circle (1.6pt); \fill (T1 -| O) circle (1.6pt);
  \draw (O) to[short, *-o] ++(0.6,0) node[right, tens] {$\Vout$};
  \draw (A.+) -- ++(-0.5,0) -- ++(0,-0.4) node[ground] {};
  \fill (N) circle (1.6pt);
\end{circuitikz}\\[2pt]
{\small (b)}
\end{minipage}
\caption{(a) Amplificador inversor generalizado (figura~4a del guion): $Z_1(s)$ y $Z_2(s)$
son impedancias arbitrarias. (b) Caso $Z_1=R$, $Z_2=R_F\parallel C$ (figura~5a del guion).}
\label{fig:generalizado}
\end{figure}


% -------------------------------------------------------------------------
\subsection{Herramienta: impedancia y transformada de Laplace}
% -------------------------------------------------------------------------
La ley del condensador, $i=C\,\dif v/\dif t$, es diferencial. La transformada de Laplace de una señal
$f(t)$, definida para $t\ge0$,
\begin{equation}\label{eq:L-def}
  F(s)=\Lt\{f\}(s)=\int_{0}^{\infty}f(t)\,\eul^{-st}\,\dif t ,
\end{equation}
la convierte en algebraica. Para demostrarlo, calculamos la transformada de la derivada integrando por
partes con $u=\eul^{-st}$ y $\dif w=\frac{\dif f}{\dif t}\,\dif t$, es decir, $\dif u=-s\,\eul^{-st}\dif t$ y $w=f(t)$:
\begin{equation}\label{eq:L-partes}
  \Lt\Bigl\{\frac{\dif f}{\dif t}\Bigr\}=\int_{0}^{\infty}\frac{\dif f}{\dif t}\,\eul^{-st}\,\dif t
  =\Bigl[f(t)\,\eul^{-st}\Bigr]_{0}^{\infty}-\int_{0}^{\infty}f(t)\,\bigl(-s\,\eul^{-st}\bigr)\,\dif t .
\end{equation}
En el límite superior, $f(t)\,\eul^{-st}\to0$ cuando $t\to\infty$ (las señales del circuito están acotadas y
$\Real s>0$); en el inferior, $\eul^{0}=1$. El término de contorno vale, por tanto, $0-f(0)$, y la integral
restante es $s$ veces \eqref{eq:L-def}:
\begin{equation}\label{eq:L-deriv0}
  \Lt\Bigl\{\frac{\dif f}{\dif t}\Bigr\}=s\int_{0}^{\infty}f(t)\,\eul^{-st}\,\dif t-f(0)=sF(s)-f(0) .
\end{equation}
Con condición inicial nula, $f(0)=0$ (condensador descargado), derivar equivale a multiplicar por $s$.
Aplicando la transformada, que es lineal, a los dos miembros de $i=C\,\dif v/\dif t$ con $f=v$:
\begin{equation}\label{eq:L-deriv}
  I(s)=C\,\Lt\Bigl\{\frac{\dif v}{\dif t}\Bigr\}=C\,\bigl[sV(s)-v(0)\bigr]
  \overset{v(0)=0}{=}sC\,V(s) .
\end{equation}
Cada elemento cumple así una «ley de Ohm» $V(s)=Z(s)\,I(s)$, con impedancias
\begin{equation}\label{eq:L-Z}
  Z_R=R,\qquad Z_C(s)=\frac{1}{sC},
\end{equation}
que se asocian como resistencias. Por linealidad, Kirchhoff, (R1) y (R2) valen para las transformadas, y en
régimen sinusoidal de pulsación $\omega=2\pi f$ basta tomar $s=\jj\omega$: $\abs{H(\jj\omega)}$ es el cociente de
amplitudes $\Vout/\Vin$ y $\arg H(\jj\omega)$ el desfase.

% -------------------------------------------------------------------------
\subsection{Ganancia del inversor generalizado}
% -------------------------------------------------------------------------
\begin{resolucion}
Se repite el cálculo del inversor con $R_k\to Z_k(s)$ (sentidos de la \figpanel{fig:generalizado}{a}).

\textbf{1. Tierra virtual} (R2): la entrada $+$ está a tierra,
\begin{equation}\label{eq:g-vm}
  V_-(s)=V_+(s)=0 .
\end{equation}

\textbf{2. Ohm en $Z_1$ y Kirchhoff en el nudo inversor}, con $I_-=0$ (R1):
\begin{equation}\label{eq:g-I}
  I_2(s)=I_1(s)=\frac{\Vin(s)-V_-(s)}{Z_1(s)}=\frac{\Vin(s)}{Z_1(s)} .
\end{equation}

\textbf{3. Ohm en $Z_2$}, del nudo inversor a la salida:
\begin{equation}\label{eq:g-vout}
  V_-(s)-\Vout(s)=Z_2(s)\,I_2(s)\quad\Longrightarrow\quad\Vout(s)=-\frac{Z_2(s)}{Z_1(s)}\,\Vin(s) .
\end{equation}

\textbf{Resultado.}
\begin{equation}\label{eq:g-H}
  \resultado{H(s)=\frac{\Vout(s)}{\Vin(s)}=-\frac{Z_2(s)}{Z_1(s)}},\qquad
  A_V(\jj\omega)=-\frac{Z_2(\jj\omega)}{Z_1(\jj\omega)} .
\end{equation}
Si $Z_1=R_1$ y $Z_2=R_2$, $A_V=-R_2/R_1$ es real y no depende de $\omega$; si hay un condensador,
$Z_C=1/(\jj\omega C)$ introduce la dependencia en frecuencia.
\end{resolucion}

% -------------------------------------------------------------------------
\subsection{Caso $Z_1=R$, $Z_2=R_F\parallel C$: filtro paso bajo}
% -------------------------------------------------------------------------
\begin{resolucion}
\textbf{1. Impedancias.} Por \eqref{eq:L-Z}, sumando las inversas del paralelo,
\begin{equation}\label{eq:f-Z2}
  \frac{1}{Z_2}=\frac{1}{R_F}+sC=\frac{1+sR_FC}{R_F}\quad\Longrightarrow\quad
  Z_1=R,\qquad Z_2(s)=\frac{R_F}{1+sR_FC} .
\end{equation}

\textbf{2. Función de transferencia.} Sustituyendo \eqref{eq:f-Z2} en \eqref{eq:g-H}:
\begin{equation}\label{eq:f-H}
  \resultado{H(s)=-\frac{R_F}{R}\,\frac{1}{1+sR_FC}},\qquad
  f_c=\frac{1}{2\pi R_FC} .
\end{equation}

\textbf{3. Módulo de la ganancia.} Con $s=\jj\omega=\jj2\pi f$ y $\abs{1+\jj x}=\sqrt{1+x^{2}}$,
\begin{equation}\label{eq:f-mod}
  \abs{A_V(f)}=\frac{R_F/R}{\sqrt{1+(f/f_c)^{2}}},\qquad
  \arg A_V(f)=180^{\circ}-\arctg\frac{f}{f_c} .
\end{equation}

\textbf{4. Límites.} De \eqref{eq:f-mod}:
\begin{align}
  f\ll f_c:&\quad \abs{A_V}\simeq\frac{R_F}{R}\quad\text{(constante; $C$ actúa como circuito abierto)},\label{eq:f-baja}\\
  f\gg f_c:&\quad \abs{A_V}\simeq\frac{R_F}{R}\,\frac{f_c}{f}=\frac{1}{2\pi fRC}\quad\text{(decrece como $1/f$; $C$ cortocircuita a $R_F$)}.\label{eq:f-alta}
\end{align}
En decibelios, por \eqref{eq:f-alta}, cada década de frecuencia resta
$20\log_{10}10=20\dB$; y en $f=f_c$, $\abs{A_V}=(R_F/R)/\sqrt2$, es decir, $3\dB$ por debajo de la meseta.

\textbf{Conclusión.} La ganancia es constante a baja frecuencia \eqref{eq:f-baja} y disminuye a
$-20\dB$/década a alta frecuencia \eqref{eq:f-alta}, con la transición en $f_c$: el circuito es un
\emph{filtro paso bajo de primer orden}, como en la figura~4b del guion (\cref{fig:bode-filtro}).
\end{resolucion}

% Figura: Bode del filtro paso bajo activo, con las frecuencias características en forma simbólica.
\begin{figure}[ht]
\centering
\begin{tikzpicture}
  \begin{semilogxaxis}[tema, name=mod, width=0.48\textwidth, height=5.4cm,
      xmin=10, xmax=1e6, ymin=-12, ymax=48, xtick={1001,100097}, xticklabels={$f_c$,$f_1$},
      ytick={0,37,40}, yticklabels={$0$,,$20\log\frac{R_F}{R}$},
      xlabel={$f$ (escala logarítmica)}, ylabel={$\abs{A_V}$ (dB)}, log basis x=10,
      title style={font=\small}, title={(a) Módulo}, clip=false]
    \addplot[asintota, domain=10:1001, samples=2] {40};
    \addplot[asintota, domain=1001:100097, samples=2] {20*log10(100097/x)};
    \addplot[azulacero, line width=1.1pt, domain=10:1e6, samples=300]
      {40-10*log10(1+(x/1001)^2)};
    \addplot[only marks, mark=*, mark size=1.6pt, rojoteo] coordinates {(1001,37)};
    \node[font=\scriptsize, anchor=north east, text=rojoteo] at (axis cs:950,35) {$-3\dB$};
    \node[font=\scriptsize, anchor=east, text=grispizarra] at (axis cs:2500,12) {$-20\dB$/década $\rightarrow$};
  \end{semilogxaxis}
  \begin{semilogxaxis}[tema, at={(mod.east)}, anchor=west, xshift=1.6cm,
      width=0.48\textwidth, height=5.4cm,
      xmin=10, xmax=1e6, ymin=80, ymax=190, ytick={90,135,180},
      xtick={1001}, xticklabels={$f_c$},
      xlabel={$f$ (escala logarítmica)}, ylabel={$\arg A_V$ (grados)}, log basis x=10,
      title style={font=\small}, title={(b) Fase}]
    \addplot[azulacero, line width=1.1pt, domain=10:1e6, samples=300] {180-atan(x/1001)};
    \addplot[only marks, mark=*, mark size=1.6pt, rojoteo] coordinates {(1001,135)};
  \end{semilogxaxis}
\end{tikzpicture}
\caption{Diagrama de Bode del filtro de la \figpanel{fig:generalizado}{b}. (a) Módulo, según
\eqref{eq:f-dB}: meseta $20\log_{10}(R_F/R)$, corte a $-3\dB$ en $f_c=1/(2\pi R_FC)$ y caída de
$-20\dB$/década (asíntota discontinua, que corta $0\dB$ en $f_1=1/(2\pi RC)$). (b) Fase, según
\eqref{eq:f-mod}: de $180^{\circ}$ a $90^{\circ}$, con $135^{\circ}$ en $f_c$.}
\label{fig:bode-filtro}
\end{figure}



\end{document}
